\chapter{파이썬에서 쓰기: \texorpdfstring{\code{christopher}}{christopher} 패키지}\label{ch:15-python-api}
\begingroup\setlength{\emergencystretch}{3em}% long identifiers: allow looser lines in this chapter only

앞의 다섯 장은 Rust 코어를 읽었다. 실제로 이 코어를 가장 많이 부르는 것은 Python이다. 노트북, 메뉴 카드
생성기(\ref{ch:32-menu-cards}장), FPGA 호스트(\ref{ch:31-fpga-host}장), 회로 시험의 기준 벡터가 모두 Python에서
코어를 부른다. 이 장은 그 사이의 두 층, Rust 함수를 Python에 내놓는 \code{crates/christopher-py}와 조작자의
언어로 감싸는 \code{python/christopher}를 읽는다.

\section{왜 두 언어인가}\label{sec:py-why}

\Why 두 언어는 서로 잘하는 것이 다르다. Rust는 컴파일해서 돌리므로 빠르고, 자료형을 컴파일 때 검사하므로 큰
계산의 기반이 되기에 알맞다. M4의 위치 6억 1천만 곳을 도는 일(\ref{sec:s-count}절)은 이쪽이어야 한다. Python은
한 줄씩 해 보며 결과를 바로 보는 데 알맞고, 그림을 그리거나 파일을 다루거나 하드웨어와 이야기하는 라이브러리가
많다. 그래서 계산은 Rust에 한 벌만 두고, 그것을 쓰는 일은 Python에서 한다.

\General 과학 계산에서 흔한 구조다. 수치 계산 라이브러리 NumPy는 겉은 Python이지만 행렬 곱셈은 C나 포트란으로
짠 코드가 한다. 사용자는 Python의 편리함을 누리고, 시간이 걸리는 부분은 컴파일된 코드가 맡는다.

\Ours 더 중요한 이유는 \emph{기준이 하나}라는 것이다. 에니그마의 배선표를 파일 하나에만 두는 것(\ref{ch:rust-enigma}장)과
같은 생각이다. 메뉴 카드가 예상하는 정지, FPGA 엔진과 대조할 정지, 회로 시험의 답안지가 모두 같은 Rust 코드에서
나온다. Python으로 봄베를 다시 짰다면 두 구현이 어긋날 수 있고, 그러면 무엇이 기준인지 말할 수 없게 된다.
그림~\ref{fig:g-py-layers}가 층의 구조다.

\begin{figure}[ht]
\centering
\begin{tikzpicture}[x=1cm,y=1cm]
  \node[blkg,minimum width=112mm,minimum height=9mm,font=\footnotesize] (core) at (0,0) {Rust 코어 \code{christopher-core}: \code{enigma}, \code{menu}, \code{bombe}, \code{check}, \code{search}};
  \node[blkn,minimum width=112mm,minimum height=9mm,font=\footnotesize] (bind) at (0,1.3) {바인딩 \code{christopher-py} (PyO3) $\to$ 확장 모듈 \code{christopher.\_core}: 함수 10개, 클래스 1개};
  \node[blkb,minimum width=112mm,minimum height=9mm,font=\footnotesize] (py) at (0,2.6) {Python 패키지 \code{christopher}: \code{Enigma}, \code{Menu}, \code{Bombe}, \code{Stop}, \code{Solution}, \code{search}};
  \foreach \x/\lab in {-4.9/노트북, -2.45/메뉴 카드, 0/FPGA 호스트, 2.45/RTL 벡터, 4.9/시험} {
    \node[blko,minimum width=21mm,minimum height=8mm,font=\footnotesize] at (\x,4.0) {\lab};
    \draw[->] (\x,3.6) -- (\x,3.05);
  }
  \draw[<->] (bind) -- (core);
  \draw[<->] (py) -- (bind);
  \node[note,anchor=west] at (0.15,0.65) {Rust 자료형};
  \node[note,anchor=west] at (0.15,1.95) {튜플, 정수, 글자열};
\end{tikzpicture}
\caption{코어 하나를 여러 곳이 쓴다. 가운데 층은 일부러 단순하게 두었다.}
\label{fig:g-py-layers}
\end{figure}

\section{경계는 튜플로}\label{sec:py-wire}

PyO3는 Rust 함수에 표시를 붙여 Python에서 부를 수 있게 해 주는 라이브러리다. \code{\#[pyfunction]}을 붙인 함수는
Python의 함수가 되고, \code{\#[pyclass]}를 붙인 구조체는 Python의 클래스가 된다. 인자와 반환값은 PyO3가 자동으로
바꿔 준다. Rust의 \code{u8}과 \code{usize}는 Python의 \code{int}로, \code{String}은 \code{str}로, \code{Vec}은
\code{list}로, 튜플은 튜플로, \code{Option}은 값 또는 \code{None}으로.

바인딩 파일의 머리 주석이 설계를 한 문단으로 밝힌다.

\begin{lstlisting}[language=Rust,caption={\code{christopher-py/src/lib.rs}: 머리 주석과 경계의 자료형.},label={lst:py-wire}]
//! The binding layer is deliberately dumb: letters cross the boundary as `str` (or `u8` 0..26 inside
//! tuples) and every structure is a plain tuple. `python/christopher` turns those into the
//! operator-facing classes. Only the Enigma is a stateful class. Long computations (bombe runs,
//! search, checking) release the GIL with `Python::detach`.

type MenuT = (usize, String, Vec<(u8, u8, u16)>, u8, u8);
type OrderT = (String, Option<String>, [String; 3]);
type PosT = (u8, u8, u8, u8);
type StopT = (PosT, u32);
\end{lstlisting}

메뉴는 (크립 위치, 크립, 선의 목록, 테스트 글자, 입력 선)의 튜플이고 선은 (평문, 암호문, 크립 위치)다. 로터 순서는
(반사판, 그리스 로터 또는 없음, 로터 셋의 이름), 드럼 위치는 코어 위치 네 개, 정지는 (위치, 테스트 레지스터)다.

\Why 경계를 튜플로 두는 이유는 두 가지다. 첫째, PyO3가 튜플과 정수를 그대로 바꿔 주므로 바인딩에 손으로 짠 변환이
거의 없다. 둘째, 사용자가 만지는 클래스(\code{Menu}, \code{Stop})는 순수 Python이므로, 출력 형식을 고치거나 편의
기능을 더할 때 Rust를 다시 컴파일할 필요가 없다. 상태를 가진 것은 에니그마 하나뿐이다. 에니그마는 글자를 칠
때마다 로터가 움직이므로 그 상태를 Rust 쪽 객체가 들고 있어야 한다. 봄베 구동은 메뉴와 로터 순서를 받아 정지를
돌려주는 함수여서 상태가 필요 없다.

\paragraph{경계에서 검사한다.} Python에서는 무엇이든 넘어올 수 있다. 정수 30을 글자라고 넘길 수도 있다. 그래서
튜플을 Rust 구조체로 바꾸는 곳에서 값을 검사한다.

\begin{lstlisting}[language=Rust,caption={\code{christopher-py/src/lib.rs}: 메뉴 튜플을 받을 때의 검사.},label={lst:py-menu-from}]
fn menu_from_py(m: &MenuT) -> PyResult<Menu> {
    let crib = alphabet::parse_letters(&m.1).map_err(err)?;
    let mut edges = Vec::with_capacity(m.2.len());
    for &(plain, cipher, pos) in &m.2 {
        if plain >= 26 || cipher >= 26 || plain == cipher {
            return Err(PyValueError::new_err(format!("invalid menu edge ({plain}, {cipher})")));
        }
        edges.push(Edge { plain, cipher, pos });
    }
    ...
}

fn err(e: Error) -> PyErr {
    PyValueError::new_err(e.to_string())
}
\end{lstlisting}

Rust 코어 안에서는 글자가 0--25라는 것을 믿고 배열의 칸 번호로 쓴다. 그 믿음이 깨지면 배열 밖을 읽게 된다(Rust는
이때 프로그램을 멈춘다). 그래서 바깥에서 들어오는 값은 문 앞에서 한 번 걸러야 한다. \code{plain == cipher}를 막는
것은 에니그마가 글자를 자기 자신으로 암호화하지 않기 때문이다(\ref{ch:rust-menu}장). 그런 선은 크립에서 나올 수 없다.

\Ours 오류는 모두 Python의 \code{ValueError} 하나로 나온다. 코어의 오류 두 가지(해석 오류 \code{Parse}와 설정
오류 \code{Config})는 \code{err}가 바꾸고, Python 쪽에서 글자를 해석하다 생긴 오류도 \code{ValueError}다. 사용자는
``설정이 잘못되었다''를 한 종류의 예외로 잡으면 된다. 같은 로터를 두 번 쓰거나, Enigma~I에 로터 VI를 끼우거나,
링에 27을 쓰거나, 플러그보드에 \texttt{AB AC}를 쓰면 모두 이 예외가 나는지를 시험
\code{test\_bad\_settings\_\allowbreak raise\_value\_error}가 본다.

\paragraph{계산하는 동안 GIL을 놓는다.} Python 인터프리터에는 GIL(global interpreter lock)이라는 자물쇠가 있어서,
한 순간에 Python 코드를 실행하는 스레드는 하나뿐이다. 봄베 구동처럼 몇 초에서 몇 분 걸리는 Rust 함수가 이
자물쇠를 쥔 채 돌면, 그동안 같은 프로그램의 다른 Python 스레드(예를 들어 화면을 그리는 스레드)가 모두 멈춘다.

\begin{lstlisting}[language=Rust,numbers=none]
let results = py.detach(|| search::run_orders(&menu, &orders, &opts, parallel));
\end{lstlisting}

\code{py.detach}는 ``이 계산 동안 자물쇠를 놓는다''는 뜻이다. 놓을 수 있는 것은 계산이 Python 객체를 하나도
만지지 않기 때문이다. 입력은 이미 Rust 구조체로 바꿔 두었고, 결과는 계산이 끝난 뒤 자물쇠를 다시 쥐고 튜플로
바꾼다. 계산 안에서 코어를 여럿 쓰는 것은 Rust의 병렬 처리(\ref{ch:14-rust-search}장)가 하므로 GIL과 무관하다.

\section{조작자의 말로: \texorpdfstring{\code{python/christopher}}{python/christopher}}\label{sec:py-classes}

패키지의 클래스는 실제 조작 순서를 따른다. 패키지 문서의 첫 문단이 그 순서다. ``크립으로 \code{Menu}를 만들고,
\code{Bombe}에 꽂고(\code{plug\_up}), 로터 순서를 돌리고(\code{run}), \code{Stop}의 목록을 읽고, 정지마다 체킹해
\code{Solution}을 얻는다. \code{search}는 이 모두를 한 번에 한다.''

\begin{figure}[ht]
\centering
\begin{tikzpicture}[x=1cm,y=1cm,
    call/.style={font=\scriptsize\ttfamily,text=cgray,midway,above}]
  \node[blkn,minimum width=19mm] (ct) at (0,0) {암호문\\크립, 위치};
  \node[blkb,minimum width=19mm] (menu) at (3.0,0) {\code{Menu}};
  \node[blkb,minimum width=19mm] (bombe) at (6.0,0) {\code{Bombe}};
  \node[blko,minimum width=19mm] (stop) at (9.0,0) {\code{Stop}\\{\scriptsize 목록}};
  \node[blkg,minimum width=19mm] (sol) at (12.0,0) {\code{Solution}};
  \draw[->] (ct) -- node[call] {from\_crib} (menu);
  \draw[->] (menu) -- node[call] {plug\_up} (bombe);
  \draw[->] (bombe) -- node[call] {run} (stop);
  \draw[->] (stop) -- node[call] {check} (sol);
  \node[note,anchor=north,align=center] at (3.0,-0.65) {\code{select}, \code{subset}\\\code{with\_test}};
  \draw[->,cblue,dashed] (ct.south) .. controls (3,-1.6) and (9,-1.6) .. node[below,note,text=cblue] {\code{search}: 메뉴부터 체킹까지 한 번에} (sol.south);
  \node[note,anchor=north] at (9.0,-0.65) {\code{trace}};
  \node[note,anchor=north] at (12.0,-0.65) {\code{enigma()}};
\end{tikzpicture}
\caption{패키지의 흐름. 점선은 같은 일을 한 번에 하는 \code{search}다.}
\label{fig:g-py-flow}
\end{figure}

몇 가지 설계가 눈에 띈다.

\begin{itemize}
\item \textbf{정지는 자기 봄베를 안다.} \code{Stop}은 만든 \code{Bombe}를 숨은 필드로 들고 있어서
\code{stop.check()}와 \code{stop.trace()}를 바로 부를 수 있다. 조작자가 정지 전표 한 장을 들고 체킹 머신으로
가던 모습이다. 정지 전표에는 메뉴를 다시 적지 않았다.
\item \textbf{드럼 위치는 글자로 보인다.} 코어는 위치를 (그리스, 왼쪽, 가운데, 오른쪽)의 수로 주지만,
\code{Stop.position}은 \texttt{"ETD"} 같은 글자다(M4이면 네 글자). 정지 전표에 적던 것이 드럼의 글자였다.
수는 숨은 필드 \code{\_pos}에 남겨 두어 다시 코어에 넘길 때 쓴다.
\item \textbf{스테커 후보의 규칙이 한곳에 있다.} Python의 \code{Stop.stecker\_\allowbreak candidates}는 Rust 쪽의
같은 이름의 메서드와 같은 규칙이다. 통전된 선이 하나뿐이면 그 선, 아니면 꺼진 선 전부다.
\item \textbf{해는 바로 쓸 수 있다.} \code{Solution.enigma()}는 복원한 키로 맞춘 \code{Enigma}를 돌려준다.
노트북은 이것으로 암호문을 처음부터 복호화해 원래 평문과 같은지 확인한다.
\end{itemize}

\paragraph{설정값을 받는 법.} 에니그마의 링 설정과 창 글자는 키 표에 글자로도, 숫자로도 적혀 있었다. 그래서
\code{\_text.letters}는 \texttt{"BUL"}, \texttt{"B U L"}, \texttt{"02 21 12"}, \code{[2, 21, 12]},
\code{["B", "U", "L"]}을 모두 같은 \texttt{"BUL"}로 바꾼다. 숫자는 기계의 눈금처럼 1부터 센다(\texttt{A}=1).
안쪽에서는 \texttt{A}=0이므로 바꾸는 곳은 \code{letter}의 한 줄 \code{return value - 1}뿐이다. 사람이 보는
번호와 계산이 쓰는 번호를 경계에서 한 번만 바꾸면, 안쪽 코드는 한 가지 약속만 알면 된다.

\section{부분 메뉴, 테스트 글자, 조각 구동}\label{sec:py-exhibit}

전시용으로 더한 세 기능은 Python에서 이렇게 보인다.

\begin{lstlisting}[language=Python,caption={\code{menu.py}: 부분 메뉴와 테스트 글자 바꾸기.},label={lst:py-subset}]
def subset(self, positions) -> "Menu":
    """Sub-menu of the edges at the given crib positions (0-based, as in :attr:`Edge.pos`).
    Keeps the crib and offset, so stops found with it can still be checked."""
    return Menu(_core.menu_subset(self._raw, [int(k) for k in positions]), self.ciphertext)

def with_test(self, test_letter: str, input_wire: str | None = None) -> "Menu":
    """Same menu with another test letter (it must be a menu letter) and input wire."""
    t = letter(test_letter)
    if chr(A + t) not in self.letters:
        raise ValueError(f"test letter {test_letter!r} is not in the menu ({self.letters})")
    offset, crib, edges, _, wire = self._raw
    w = wire if input_wire is None else letter(input_wire)
    return Menu((offset, crib, edges, t, w), self.ciphertext)
\end{lstlisting}

\code{subset}은 코어의 \code{Menu::subset}(\ref{sec:s-subset}절)을 그대로 부른다. 새 \code{Menu}에 암호문을
넘겨주는 것을 눈여겨볼 만하다. 체킹에는 암호문 전체가 필요하므로 메뉴가 그것을 함께 들고 다닌다.

\code{with\_test}에 해당하는 Rust 함수는 없다. 메뉴의 튜플에서 테스트 글자와 입력 선 자리만 바꿔 새 튜플을 만들면
되기 때문이다. 검사는 하나다. 테스트 글자는 메뉴에 있는 글자여야 한다. 메뉴에 없는 글자의 케이블에는 스크램블러가
하나도 이어져 있지 않다. 그런 케이블을 테스트 레지스터에 꽂는 것은 메뉴를 잘못 읽은 것이므로 여기서 막는다.

\Why 테스트 글자를 바꾸는 이유는 전시 카드에 있다. 카드는 ``테스트 글자 \texttt{D}의 케이블을 테스트 레지스터에
꽂는다''처럼 꽂을 자리를 인쇄해 둔다(\ref{ch:32-menu-cards}장). 소프트웨어가 예상한 정지와 실물이 낸 정지를
비교하려면, 소프트웨어도 정확히 그 글자와 그 입력 선으로 돌아야 한다. 원리로는 메뉴의 어느 글자를 테스트 글자로
삼아도 맞는 위치에서 정지하지만(\ref{sec:propagate}절의 논증은 테스트 글자가 어느 것이든 성립한다), 거짓 정지는
글자에 따라 달라질 수 있다. 시험
\code{test\_with\_test\_letter}는 글자가 바뀌고 선은 그대로인지, 메뉴에 없는 글자를 주면 \code{ValueError}가 나는지
본다.

\paragraph{조각 구동.} \code{Bombe.run}과 \code{search}는 \code{fixed\_right}를 받는다(\ref{sec:s-slice}절).

\begin{lstlisting}[language=Python,caption={\code{bombe.py}: 느린 드럼의 위치를 글자나 수로 받는다.},label={lst:py-fixed}]
def _core_offset(value: str | int | None) -> int | None:
    """A drum core given as a letter ("X") or an offset (23)."""
    if value is None:
        return None
    if isinstance(value, str):
        return letter(value)
    if not 0 <= int(value) < 26:
        raise ValueError(f"drum offset must be 0..25, got {value}")
    return int(value)
\end{lstlisting}

글자 \texttt{"X"}로도, 코어 위치 23으로도 줄 수 있다. 여기서 수는 1부터 세는 눈금이 아니라 0부터 세는 코어
위치다. 정지의 드럼 글자를 그대로 넘기는 것이 가장 흔한 쓰임이어서, 메뉴 카드 생성기는 정답 위치의 셋째 글자를
그대로 넘긴다. 바인딩도 26 이상이면 한 번 더 거른다(\code{check\_right}). 시험
\code{test\_slice\_run\_fixed\_right\_drum}은 \ref{sec:s-slice}절의 Rust 시험과 같은 것을 Python에서 본다.

\begin{lstlisting}[language=Python,caption={\code{tests/test\_python\_api.py}: 조각 구동(일부).},label={lst:py-slice-test}]
full = bombe.run("B III-V-II")
true = [s for s in full if s.check()]
right = true[0].position[-1]
sliced = bombe.run("B III-V-II", fixed_right=right)
assert bombe.positions_tested == 676
assert [s.position for s in sliced] == [s.position for s in full if s.position[-1] == right]
\end{lstlisting}

전체 구동의 정지 가운데 체킹을 통과하는 것(참인 정지)을 찾고, 그 드럼 글자의 마지막 글자(오른쪽 드럼)로 조각을
돌린다. 시험한 위치가 676이고, 조각의 정지가 전체 정지 가운데 오른쪽 드럼이 같은 것과 순서까지 같아야 한다.

\section{회로 시험의 답안지: \texorpdfstring{\code{export\_vectors}}{export\_vectors}}\label{sec:py-vectors}

회로를 시험하려면 ``이 위치에서 676가닥이 어떻게 되어야 하는가''라는 답이 필요하다. \code{vectors.py}의
\code{export\_vectors}가 그 답을 파일로 쓴다. 한 줄에 위치 하나씩, JSON 객체 하나다.

\begin{lstlisting}[numbers=none]
{"wheel_order": "B II-IV-V", "position": [g, l, m, r], "edges": [[plain, cipher, pos], ...],
 "test_letter": t, "input_wire": w, "live": [26 ints], "stop": bool}
\end{lstlisting}

\code{live}는 \code{Bombe.trace}의 결과, 곧 조기 종료 없이 끝까지 퍼뜨린 676비트다. 위치를 따로 주지 않으면
그 로터 순서의 정지 전부와 정지가 아닌 위치 64곳을 무작위로 골라 쓴다.

\Why 정지가 아닌 위치도 넣는 데에는 이유가 있다. 회로가 늘 ``정지''라고 답하는 고장이 있다고 하자. 정지인
위치만 시험하면 이 고장을 잡지 못한다. 시험의 답안지에는 ``예''와 ``아니오''가 둘 다 있어야 한다. 무작위로 고르되
씨앗(\code{seed})을 고정해 두므로, 같은 인자로 다시 만들면 같은 파일이 나온다. 파일이 바뀌면 코드가 바뀐 것이다.

\General 계산기를 시험할 때 $2+2$만 넣어 보지 않는다. 답이 0인 것, 음수인 것, 자리올림이 있는 것을 섞는다. 같은
시험지를 내년에도 쓸 수 있도록 문제 목록은 적어 둔다.

\Ours 이 파일들이 \code{data/vectors/rtl/}에 있고(\code{rtl/tb/export\_vectors.py}가 같은 형식을 넓혀 만든다),
기준 RTL의 \code{bombe\_core} 시험과 FPGA 엔진의 스냅숏 시험(\ref{ch:31-fpga-host}장)이 이것과 대조한다. 시험
\code{test\_export\_vectors}는 \code{sample=5}로 파일을 써서, 줄 수가 정지 수 $+5$이고 \code{"stop"}이 참인 줄이
정확히 정지 수만큼인지 본다.

\section{노트북으로 따라가기}\label{sec:py-notebook}

\code{notebooks/01\_crib\_to\_plaintext.ipynb}는 이 패키지로 봄베 조작의 전 과정을 따라 하는 예제다. 실행 결과가
노트북에 저장되어 있어서 열기만 해도 읽을 수 있다. 각 절이 부르는 API와 확인하는 것을 표로 정리한다. 결과의
해석(정지 157회, 링 설정이 \texttt{AAC}로 나오는 이유 등)은 \ref{ch:rust-menu}--\ref{ch:rust-check}장에서
다루었으므로 되풀이하지 않는다.

\begin{center}
\small
\begin{tabularx}{\linewidth}{>{\raggedright\arraybackslash}p{30mm}>{\raggedright\arraybackslash}p{40mm}>{\raggedright\arraybackslash}X}
\toprule
노트북의 절 & 부르는 것 & 확인하는 것 \\
\midrule
1. 에니그마 & \code{Enigma(...).encrypt}, \code{reset}, \code{positions} & 상호성(같은 키로 다시 치면 평문), 121글자 중 같은 자리 같은 글자 0개, 이중 스테핑 \texttt{ADU}$\to$\texttt{BFX} \\
2. 크립 정렬 & \code{possible\_offsets} & 106곳 중 60곳이 남음, 실제 위치 28이 그 안에 있음 \\
3. 메뉴 & \code{Menu.from\_crib}, \code{select(14)}, \code{components} & 선 16개의 메뉴, 떨어진 덩어리 \texttt{FO}, \texttt{NS}, 다듬은 메뉴(선 14, 루프 1) \\
4. 봄베 & \code{Bombe.plug\_up}, \code{run("all")}, \code{trace(matrix=True)} & 120 순서, 2,109,120 위치, 정지 목록, 대각선 보드를 끄면 17,576곳 전부 정지, 676가닥의 그림 \\
5. 체킹 & \code{Stop.check}, \code{Solution.enigma} & 정지 하나만 통과, 복원한 키로 전문 전체 복호화 \\
6. 넘어감 & \code{search}(넘어감 가설 켬) & 가설 없이 정지 0회, 가설로 해 1개 \\
7. M4 & \code{search(model="m4", ...)}, 1945년 실제 전문 & 로터 순서 6개로 키 복원, 실제 U보트 전문의 키 전체 복원 \\
\bottomrule
\end{tabularx}
\end{center}

\Ours 노트북은 결과를 눈으로 보여 주는 것에 그치지 않는다. 중요한 곳마다 \code{assert}가 있다. 상호성
(\code{enigma.reset().decrypt(CIPHER) == PLAIN}), 복원한 키의 복호문(\code{found.enigma().decrypt(CIPHER) ==
PLAIN})이 그렇다. 노트북을 다시 실행해서 끝까지 가면 그 자체로 시험 하나를 통과한 것이다. 시간이 찍힌
출력(``0.32초에 탐색'' 등)은 실행한 컴퓨터에 따라 다르다.

\begin{keybox}[title={이 장의 정리}]
\begin{itemize}
\item 계산은 Rust에 한 벌, 쓰는 일은 Python에서. 기준이 하나여야 여러 구현을 대조할 수 있다.
\item 경계는 튜플과 정수로 단순하게 두고, 들어오는 값은 경계에서 검사한다. 오류는 \code{ValueError} 하나로.
\item 긴 계산은 GIL을 놓고 돈다. 계산이 Python 객체를 만지지 않으므로 놓을 수 있다.
\item 사람이 보는 표기(1부터 세는 숫자, 드럼 글자)와 계산의 표기(0부터 세는 수)는 경계에서 한 번만 바꾼다.
\item 시험의 답안지에는 정지인 위치와 아닌 위치가 둘 다 들어간다.
\end{itemize}
\end{keybox}
\endgroup
